\cotefile{102}
\scanpage{2}
\hfill\uline{Recollements}

% archive sheet number: 1

\scanpage{3}
Pour EGA VI

\uline{Recollements d'objets d'un site}.

Soit $S$ un site\inserted{ standard (i.e. à topologie moins fine que la canonique}, muni d'une partie $\underline{M}$ de l'ensemble des flèches, satisfaisant les conditions habituelles :

a) les flèches de $M$ sont quarrab\deleted{k}\handwritten{l}es et stables par changement de base,

b) $M$ contient les flèches\handwritten{\ensuremath{\vert}}identiques et est stable par composition,

c) une flèche $u:X\to Y$ appartient à $M$ si elle y appartient localement sur $Y$,

et en plus une condition de saturation:

d) Soit $u_{\handwritten{i}}:X_{\handwritten{i}}\to\deleted{Y}\handwritten{X}$ un\handwritten{e}\inserted{\handwritten{ famille couvrante de}} \deleted{$Y$-}morphisme\handwritten{s} dans $\deleted{S}\handwritten{M}$, \deleted{supposons \handwritten{\unsure{les}} $X_{\handwritten{i}}\unsure{\to|\,} \unsure{Y|X}$ et \unsure{les|dans}}\inserted{\handwritten{(i.e, par a) et c), les}} $X_{\handwritten{i}}\times_{\deleted{Y}\handwritten{X}}X_{\handwritten{j}}\xrightarrow{\mathrm{pr}_{1}}X_{\handwritten{i}}$ dans $M$\handwritten{)}, \inserted{\handwritten{et supposons les $uu_{\unsure{i}}:X_i\to Y$ dans $M$,}} alors $u\inserted{\handwritten{{}:X\to Y}}$ est dans $M$.

% slanted
\margin{\handwritten{essayer de ne pas le mettre ici, en saturant plus tard en $\widetilde M$, et formulant $\tilde e$ de façon un peu différente : $u$ est dans $M$ sss on peut trouver des $u_i:X_i\to X$ dans $M$ avec $uu_i$ \unsure{ds|dans} $M$.}}

\uline{Proposition} I\handwritten{l} existe une partie $\overset{\handwritten{\sim}}{M}$ de l'ensemble des flèches du topos $\overset{\handwritten{\sim}}{S}$, déterminée de façon unique par les propriétés suivantes: $\overset{\handwritten{\sim}}{M}$ satisfait aux conditions \inserted{\handwritten{$\tilde a,\tilde b,\tilde c,\tilde d$}} analogues à a) à d), et de plus aux conditions suivantes:

$\overset{\handwritten{\sim}}{e}$) Si $u:X\to Y$ est une flèche de $S$, alors $u$ est dans $M$ sss $u$ est dans $\overset{\handwritten{\sim}}{M}$.

$\overset{\handwritten{\sim}}{f}$) Si $u:F\to X$ est une flèche dans $\overset{\handwritten{\sim}}{\mathbb{M}}$, avec $X\handwritten{\in}\deleted{S}\handwritten{\mathrm{Ob}\,S}$, alors il existe une famille couvrante $u_i:X_i\to F$ de flèches dans $\overset{\handwritten{\sim}}{M}$, avec $X_i\handwritten{\in}S$\handwritten{.} \deleted{et $uu_i\handwritten{\unsure{\in}}\overset{\handwritten{\sim}}{M}$.}

Démontrons l'unicité, et montrons que l'on peut décider sans ambiguité si une flèche $u:F\to G$ de $\overset{\handwritten{\sim}}{S}$ est dans $\overset{\handwritten{\sim}}{M}$. Pour ceci, distinguons deux cas:

1º) Cas où le morphisme diagonal $F\to F\times_G\deleted{\unsure{G|E}}F$ est représentable. \deleted{\unsure{U|X}\illegible{1}}

On sait qu'on peut trouver une famille couvrante $X_i\to F$, les $X_i\handwritten{\in} S$,

% archive sheet number: 2

\scanpage{4}
% The typed line at the head of the sheet is keyed by a pen loop to the first case below.

et en vertu de $\overset{\handwritten{\sim}}{c}$), on a $u\handwritten{\in} \overset{\handwritten{\sim}}{M}$ sss pour tout $i$, $F\times_G X_i\to X_i$ est dans $\overset{\handwritten{\sim}}{M}$. Cela nous ramène au cas où $G$ est de la forme $X\handwritten{\in} S$. Je dis qu'alors $u$ est dans $\overset{\handwritten{\sim}}{M}$ sss il existe une famille couvrante $u_i:X_i\to F$ avec \deleted{$u_i\unsure{\in} \overset{\handwritten{\sim}}{M}$, $uu_i\unsure{\in} M$,} $X_i\handwritten{\in} S$, $uu_i\handwritten{\in} M$, \deleted{\unsure{\&test|et est}} et les $X_i\times_F X_j\xrightarrow{\mathrm{pr}_1}X_i$ dans $M$ (NB cela a un sens, car l'hypothèse faite sur $u$ implique que $X_i\times_F X_j$ est représentable, comme image inverse \inserted{dans $X_i\times_X X_j$} de la diagonale de $F$ sur $X$)\deleted{dans} La condition est nécessaire:\deleted{\unsure{=|:}}on trouve les $u_i$ grâce à $\overset{\handwritten{\sim}}{f}$), alors par $\overset{\handwritten{\sim}}{a}$) la projection $X_i\times_F X_j\to X_i$ est dans $\overset{\handwritten{\sim}}{M}$, donc dans $M$ par $\overset{\handwritten{\sim}}{e}$)\handwritten{\unsure{6|$\ell$}}\inserted{\handwritten{;}les $uu_i$ sont dans $M$ par $\overset{\handwritten{\sim}}{b}$) et $\overset{\handwritten{\sim}}{e}$).} La condition est suffisant\unsure{e|é} grâce à $\overset{\handwritten{\sim}}{e}$) et $\overset{\handwritten{\sim}}{d}$).

2º) Cas général. Procédant comme plus haut, on estbramené au cas où $G$ est un $X\handwritten{\in} S$, et dans ce cas, on voit de même que $u\handwritten{\in} \overset{\handwritten{\sim}}{M}$ sss il existe une famille couvrante de flèches $u_i:X_i\handwritten{\to} F$, avec $uu_i\handwritten{\in} M$, et les $Z_{ij}=X_i\times_F X_j\xrightarrow{\mathrm{pr}_1}X_i$ dans $\overset{\handwritten{\sim}}{M}$. Or ici $Z_{ij}$ est un sous-\deleted{\unsure{foncteur}}faisceau\unsure{.} du faisceau représentable $Z'_{ij}=X_i\times_X X_j$, donc \deleted{sa}son \deleted{diagonale}\inserted{morphisme diagonal} sur $X_i$ est image inverse du morphisme diagonal pour $Z'_{ij}$ sur $X_i$, qui est représentable car le carré fibré de $Z'_{ij}$ sur $X_i$ l'est, donc delui de $Z_{ij}$ sur $X_i$ l'est aussi. On est donc ramené au cas 1º).

Pour achever de prouver la proposition, il reste à montrer que l'ensemble de flèches $\overset{\handwritten{\sim}}{M}$ construit par les critères précédents satisfait bien aux conditions $\overset{\handwritten{\sim}}{a}$) à $\overset{\handwritten{\sim}}{f}$), ce qu'on laisse au lecteur le soin de faire.

\uline{Définition} On dit qu'un objet $F$ de $\overset{\handwritten{\sim}}{S}$ peut \deleted{\unsure{a}}s\uline{'obtenir par recollement d'objets de} $S$ \uline{à l'aide de morphismes dans} $M$, si on peut trouver une famille couvrante $u_i:X_i\to F$, avec $X_i\handwritten{\in} S$ et $u_i\handwritten{\in} \overset{\handwritten{\sim}}{M}$. On désignera par $ {\overset{\handwritten{\sim}}{S}}_M$ la sous-catégorie pleine de $\overset{\handwritten{\sim}}{S}$ ayant comme objets ces $F$.

\deleted{\unsure{n}}On a alors \deleted{\unsure{le}}la propriété de stabilité suivante:

\deleted{\uline{Proposition} \quad 1) Soit $F\quad\mathrm{Ob}\,\overset{\handwritten{\sim}}{S}$, et $F_{i}\quad F$ une famille couvrante}

% archive sheet number: 3

\scanpage{5}
\uline{Proposition} \quad (i) Tout faisceau représentable $F\handwritten{\in} S$ est dans $\overset{\handwritten{\sim}}{S}_M$. \deleted{Si}

\deleted{on a}\inserted{(ii) Soit} un morphisme $u:F\handwritten{\to}G$ \deleted{de faisceaux,} dans $\overset{\handwritten{\sim}}{M}$, \deleted{et} Si $G\handwritten{\in}\mathrm{Ob}\,\overset{\handwritten{\sim}}{S}_M$, alors $F\handwritten{\in}\mathrm{Ob}\,\overset{\handwritten{\sim}}{S}_M$, \handwritten{et}\deleted{L}\handwritten{l}a réciproque est vraie si $u$ est couvrant.

(iii) Soit $u:F\to G$ un morphisme couvrant, tel que \deleted{$F\unsure{\in}\mathrm{Ob}\,\overset{\handwritten{\sim}}{S}_M$ et} $F\times_G F\xrightarrow{\handwritten{\mathrm{pr}_1}}F$ soit dans $\overset{\handwritten{\sim}}{M}$, alors $G$ est dans $\overset{\handwritten{\sim}}{S}_M$ sss $F$ l'est.

(i\deleted{ii}\handwritten{v}) Une somme $\handwritten{\coprod_i}F_i$ de faisdeaux est dans $\overset{\handwritten{\sim}}{S}_M$ sss les $F_i$ le sont.

(v) \deleted{Soit $I$ une catégorie finie, telle que}\inserted{Supposons} $S$ soit stable par \deleted{lim de} produit de deux objets (resp. par produit fibré de deux objets sur un \deleted{type $I$. Alors il en est de même de $\overset{\handwritten{\sim}}{S}_M$.} troisième, resp. par objet final, resp. par $\underset{\handwritten{\leftarrow}}{\lim}$ finies) alors il en est de même de $\overset{\handwritten{\sim}}{S}_M$.

(i) est trivial, \deleted{Prouvons} ainsi que la réciproque dans (ii). La partie directe est facile .... (iii) résulte de (ii), car $u$ étant couvrant, la condition $\mathrm{pr}_1\handwritten{\in} \overset{\handwritten{\sim}}{M}$ équivaut à $u\handwritten{\in} \overset{\handwritten{\sim}}{M}$. Prouvons (iv): si les $F_i$ sont dans $\overset{\handwritten{\sim}}{S}_M$, il en est de même de la somme: trivial. \deleted{\unsure{D}}Réciproque: l'inclusion $F_i\handwritten{\to} F=\handwritten{\coprod_i}F_i$ est dans $\overset{\handwritten{\sim}}{M}$ grâce à $\overset{\handwritten{\sim}}{c}$), donc si $X$ est représentable et $X\to F$ est dans $\overset{\handwritten{\sim}}{M}$, alors $X_i=X\times_F F_i \handwritten{\in} X$ est dans $\overset{\handwritten{\sim}}{M}$, donc $X_i$ peut se couvrir par des $X_{ij}\handwritten{\in} S$ avec $X_{ij}\to X_i$ dans $\overset{\handwritten{\sim}}{M}$ par $\overset{\handwritten{\sim}}{f}$); alors les composés $X_{ij}\to X_i\to F_i$ sont dans $\overset{\handwritten{\sim}}{M}$, d'autre part les $X$ envisagés couvrent $F$, donc les $X_{i\deleted{\unsure{i}}j}$ \deleted{\unsure{($X,i$)}} ($X,i$ variables) couvrent $F_i$, donc les $X_{ij}$ couvrent $F_i$, qui est donc dans $\overset{\handwritten{\sim}}{S}_M$. On laisse aussi (v\unsure{)|à} au rédacteur (attention, il est douteux que ça marche pareil pour un noyau, car il ne sera pas vrai nécessairement qu\deleted{\unsure{e}}'un \deleted{\unsure{homo}}morphisme de doublets dans $\overset{\handwritten{\sim}}{S}$ \inserted{ou même de $S$} induisant des morphismes $\handwritten{\in} \overset{\handwritten{\sim}}{M}$ sur les sommets induise un morphisme \deleted{\unsure{de}} $\handwritten{\in} M$ sur les noyaux des double flèches).

\uline{Corollaire} \quad $\mathrm{Ob}\,\overset{\handwritten{\sim}}{S}_M$ est le plus petit sous-ensemble de $\mathrm{Ob}\,S$ contenant $\mathrm{Ob}\,S$ et stable par \inserted{sommes et par} passage au quotient par une relation d'équivalence à $\mathrm{pr}_1$ qui est dans $\overset{\handwritten{\sim}}{M}$. Par suite, \deleted{\unsure{il|la}} $\overset{\handwritten{\sim}}{S}_M$ ne change pas si on remplace $S$ par une sous-catégorie $S'$ de $\overset{\handwritten{\sim}}{S}$ contenue dans $\overset{\handwritten{\sim}}{S}_M$, et $M$ par $M'=\overset{\handwritten{\sim}}{M}\handwritten{\cap}\mathrm{Fl}\,S'.$

% archive sheet number: 4

\scanpage{6}
\deleted{telle que les $F_i\times_F F_j\quad F_i$ soient dans $\overset{\handwritten{\sim}}{M}$. Pour que l'on ait $F\quad\mathrm{Ob}\,\overset{\handwritten{\sim}}{S}_M$ il fet s. que pour tout $i$, on ait $F_i$}

% archive sheet number: 5

\scanpage{7}
Exemples: Si $S$ est la catégorie des schémas affines, avec une topologie plus fine que Zar, $M$ l'ensemble des immersions ouvertes, \deleted{alors la const}\deleted{la construction donn\unsure{é|e}e} (ou, en l'absence de d), les immersions de la forme $\mathrm{Spec}(A_f)\handwritten{\to}\mathrm{Spec}(A)$), alors la construction donne une catégorie $\overset{\handwritten{\sim}}{S}_M$ équivalente à celle des schémas. Si on prend une topologie plus fine que la topologie étale, et pour $M$ \inserted{\deleted{\unsure{h}}l'ensemble} des morphismes étales (qui satisf\deleted{\unsure{on}}ait la condition d)), alors on trouve les ``espaces algébriques'' de Artin, en ``recollant des spectres d'anneaux \deleted{pour \unsure{l}} par des morphismes étales''. Le dernie\deleted{\unsure{f}}r corollair\unsure{e|è} nous dit qu'on aurait trouvé le même résultat en partant de la catégorie $S'$ des schémas, au lieu de celle $S$ des schémas affines. \handwritten{Dans ce cas, les morphismes de $\widetilde M$ sont dits \uline{étales} [la top. sur $S$ étant la top. étale). \deleted{\illegible{1}}}

\uline{Proposition} \quad Supposons que $M$ satisfasse la condition draconienne suivante: soient $u:X\to Y$ et $v:Y\to Z$ des flèches composables dans $S$, \deleted{\unsure{alor|alors}} alors $v\handwritten{\in} M$ et $vu\handwritten{\in} M$ implique $u\handwritten{\in} M$. Pour un objet $X$ de $S$, soit $S(X,M)$ la sous-catégorie de $S_{\handwritten{/}X}$ \deleted{formé} formée des $X'\to X$ qui sont dans $M$ (donc la condition précédente dit que tous les morphismes de $S(X,M)$ sont dans $M$); munissons-là de la topologie induite par $S$, et considérons le topos $S(X,M)^{\handwritten{\sim}}$. Alors \deleted{le foncteur $(\overset{\handwritten{\sim}}{S}_M)_{\handwritten{/}X}\quad S(X,M)^{\handwritten{\sim}}$ composé de}
\diagram{\begin{tikzcd}[ampersand replacement=\&,column sep=large]
\deleted{(\overset{\handwritten{\sim}}{S}_{M\unsure{x}})_{\handwritten{/}X}}
\arrow[r,"{\deleted{\mathrm{incl.\ can.}}}"] \&
\deleted{\overset{\handwritten{\sim}}{S}_{\handwritten{/}X}}\arrow[r,"{\deleted{\handwritten{\approx}}}"] \&
\deleted{(S_{\handwritten{/}X})^{\handwritten{\sim}}}\arrow[r,"{\deleted{\mathrm{restr}}}"] \&
\deleted{S(M,X)^{\handwritten{\sim}}}
\end{tikzcd}}
\deleted{est une équivalence de catégories.} la sous-catégorie pleine de $\overset{\handwritten{\sim}}{S}_{\handwritten{/}X}$ formée des $F\handwritten{\to}X$ qui sont dans $\overset{\handwritten{\sim}}{M}$ \inserted{\handwritten{[d'où résulte $F\in\mathrm{Ob}\,\widetilde S_M$ !]}}, est \deleted{équivalent} un \uline{topos}, équivalent par le foncteur restriction à $S(\deleted{M}X,\deleted{X}M)$ au topos $S(X,M)^{\handwritten{\sim}}$.

% archive sheet number: 6

\scanpage{9}
\uline{Questions pour Artin}

1°) $S$ \deleted{strictement} local \deleted{\unsure{noeth|noeth.}} \uuline{hens.} noeth, $S'\to S$ \deleted{\unsure{fpt.}} plat. surjectif, $S'$ noeth, $X'/S'$ \deleted{\unsure{de p.f.|loc. de t.f.}}\inserted{\uline{propre}} avec donnée de descente rel. à $S'/S$. Supposons que la donnée de descente sur $X'_s$ soit effective ($s$ pt. fermé de $S$). Alors la donnée de descente sur $X'$ est effective ?

Cas où $S'=\widehat{S}$ ?

2°) Peut-on faire la descente fpqc des \uline{variétés} \inserted{utiliser \unsure{théor.|loc.} de pts. fixes,}, p.ex. à partir d'un complété d'un anneau local ? [i.e. la propriété pour un faisceau fpqc loc. de prés. finie d'être une variété est-elle loc. \unsure{fpqc|fppc} sur $S$ ?]

% slanted
\margin{Le Pb est la comm. aux $\underset{\leftarrow}{\lim}$ (\unsure{adiques artiniennes})}

3°) $S$ comme dans \uline{1°}, $X$ une \uline{variété} \deleted{sur, \unsure{loc. de pr. f.|loc. de t.f.}}\inserted{propre} sur $S$. Si $X$ est représentable loc. \unsure{fppf|fppqc|fpqc} sur $S$, \inserted{et $X_s$ représentable,} alors $X$ est-il représentable ? Evidemment, si la réponse à 2° \deleted{\unsure{et 3°|est}} était affirmative, les questions 1° et 3° seraient équivalentes...

% slanted
\margin{\unsure{Exemple|Exemples} de $X$ plat p\unsure{rés.} propre sur $S$ local complet, $X_s$ repr. mais $X$ non repr. ?? Exemple de $X$ \unsure{vari.|variétés} \deleted{\unsure{de}} \unsure{pfppf|propres} sur $S$ noeth, à fibres repr., qui ne soit repr. ?}

4°) Descente fpqc des \inserted{(``petits'')} faisceaux étales.

Lié à la question de \unsure{critériser|conditions sur} les foncteurs $(\mathrm{Sch})_{/S}^{\circ}\to(\mathrm{Ens})$ qui proviennent de \deleted{petits} faisceaux étales sur $S$. Une réponse affirmative à 2° (\unsure{où|si|i.e.} on prendrait ``variété'' dans le sens général, sans exiger que le \inserted{morphisme} \unsure{diagonale|diagonal} soit une immersion) en impliquerait une à 4°.

% archive sheet number: 8

\scanpage{10}
\mspage{65}
\hfill n° 387

tel que $\handwritten{\lambda<}\mu$. Pour toute extension $E$ de $K$ et tout prolongement $w$ de $v$ à $E$, désignons par $D_E(\handwritten{\lambda})$ l'ensemble des $z\handwritten{\in}E$ tels que $v(z-a_{\handwritten{\lambda}})\handwritten{\geq\unsure{\gamma|\delta}_{\lambda}}$, où $\handwritten{\gamma_\lambda}$ est la valeur commune des $v(a_{\handwritten{\lambda}}-a_\mu)$ pour $\handwritten{\lambda<}\mu$ ; l'hypothèse est donc que $\handwritten{\bigcap_{\beta\in B}}D_K(\handwritten{\beta})=\handwritten{\varnothing}$ (§ 5, exerc. 5 a)). Soit $\handwritten{\Omega}$ une clôture algébrique de $K$ et soit $v_{\unsure{o|0}}$ un prolongement de $v$ à $\handwritten{\Omega}$ ; montrer que si $P$ est un polynôme de $K\handwritten{[}X\handwritten{]}$, pour qu'il existe $\handwritten{\lambda\in}B$ tel que $v(P(a_\mu))=v(P(a_{\handwritten{\lambda}}))$ pour tout $\handwritten{\lambda\leq}\mu$, il faut et il suffit qu'aucun zéro de $P$ dans $\handwritten{\Omega}$ n'appartienne à $\handwritten{\bigcap_{\beta\in B}}D_{\handwritten{\Omega}}(\handwritten{\beta})$. Si $P$ possède cette propriété et si $E$ et $w$ ont les mêmes significations que ci-dessus, montrer que pour tout $z\handwritten{\in\bigcap_{\beta\in B}}D_E(\handwritten{\beta})$, on a $w(P(z)-P(a_\mu))\handwritten{>}w(P(z))=v(P(a_\mu))$ dès que $\mu$ est assez grand (décomposer $P(X)$ en facteurs dans $\handwritten{\Omega[}X\handwritten{]}$). En déduire que l'on a une des deux situations suivantes :

1° Ou bien $\handwritten{\bigcap_{\beta\in B}}D_{\handwritten{\Omega}}(\handwritten{\beta})\handwritten{\neq\varnothing}$, et si $\handwritten{\unsure{\vartheta|\theta}}$ est un élément de cette intersection dont le degré sur $K$ est le plus petit possible, $K(\handwritten{\vartheta})$ est une extension immédiate de $K$ distincte de $K$.

2° Ou bien $\handwritten{\bigcap_{\beta\in B}}D_{\handwritten{\Omega}}(\handwritten{\beta})=\handwritten{\varnothing}$ ; il y a alors une valuation $v'$ de $K(X)$ prolongeant $v$ et telle que $v(P(X))=v(P(a_\mu))$ pour $\mu$ assez grand, et $K(X)$, muni de $v'$, est une extension immédiate de $K$ (utiliser le critère de a)).

e) Déduire de c) et d) que pour que $K$ soit linéairement compact, il faut et il suffit que $K$ ne possède aucune extension immédiate autre que $K$.

\begin{center}\texttt{-:-:-:-}\end{center}

% archive sheet number: 9

\scanpage{11}
1. Morphismes représentables. Morphismes $P$-représentables.

Abus de langage : imm., im. ouverte, imm. fermée, lisse, net, étale, \unsure{Osons} dire : repr. par ....\} Critère de représentabilité par ``recollement de morceaux''. (SGAD \uline{XI} 3.5.)

2. [Faisceaux qu. \unsure{cpct|cpt}, quasi-séparés. Morphismes qu. \unsure{cpct|cpt}, qu. sép. de faisceaux]. Morphismes \inserted{loc.} de prés. finie, de prés. finie. \margin{$[\,]=$ petits caractères}

3. Morphismes $\mathrm{form}^{\mathrm{t.}}$ \unsure{lisse|lisses}, $\mathrm{form}^{\mathrm{t.}}$ nets, $\mathrm{form}^{\mathrm{t.}}$ étales, et variantes infinitésimales. Critères standart....

4. Critères pour qu'un morphisme soit un monomorphisme, un bimorphisme, une immersion ouverte.

5. Conditions de séparation:
\[
\left\{\begin{aligned}
&\text{morph. diag. repr.}\\
&\text{morphisme diagonal repr. \inserted{par imm., par imm. qu. \unsure{cpcte|cpte}, \inserted{par imm. de prés. finie,}} par imm. fermée}\\
&\text{Relations avec \uline{critères valuatifs}}
\end{aligned}\right.
\]

% The arrow brings item 6 up from the space between 7 and 8.
\inserted{6. Conditions de propreté.}

7. Espaces annelés et foncteurs \unsure{, en|sur} \unsure{``Sch''|$(\mathrm{Sch})_{/}$}. Morphismes surjectifs, \deleted{\unsure{ensemblistes}}\deleted{radiciels} radiciels (morphisme diagonal est surjectif).

[Une immersion ouverte surjective est un isom].

8. Conditions de connexion.

Principe terminologiques pour les morphismes de faisceaux sur $(\mathrm{Sch})$ : a) Conditions locales b) Si $f:F\to G$ est représentable, cela doit cadrer avec la définition déjà acceptée pour \deleted{\unsure{sché}} les \unsure{morphismes} de schémas [$\forall\ X\to G$, le morphisme de schémas $\deleted{\unsure{F}}X\times_G F\to X$ a la propriété]

% archive sheet number: 10

\scanpage{12}
dans une extension radicielle convenable de $\inserted{k(y)=}R(Y)$, on peut supposer que $(X_y)_{\mathrm{réd}}$ est séparable sur $k(y)$, et remplaçant $X$ par $X_{\mathrm{réd}}$, on sera sous les conditions du \unsure{th|th.}, et on gagne.

La deuxième assertion se déduit de la troisième par réduction au cas noethérien. Pour la troisième, \deleted{distinguons les deux cas. \unsure{Si} l'anneau \unsure{principal|princip. idéal} de $Y$ est \unsure{fini}, \inserted{$Y$ \unsure{qu. cpt qu. sép.},} on est ramené au cas $Y$ intègre, \unsure{Soit} pt gén. $y$. Faisant une extension radicielle finie $y'$ de \unsure{$y$|$\eta$|$Y$} et \unsure{normaliant|normalisant}} \deleted{Dans le premier cas.} \inserted{On peut supposer $Y$ réduit.} \deleted{On \unsure{soit} $M=\mathrm{Hilb}$ \unsure{peut supposer}} Notant $M=\underline{\mathrm{Hilb}}_{X/Y}$, on a alors une \uline{section} \uline{ensembliste} $\varphi$ de $M$ sur $Y$, en associant à \unsure{tt|$\forall$} $y\in Y$ le point \inserted{$\varphi(y)$} de $M$, \unsure{not|sit} sur $y$ \deleted{tel que} défini par $(X_y)_{\mathrm{réd}}\subset\deleted{\unsure{X}}X_y$. \deleted{Soit D'} Soit $Y'=\varphi(Y)$. Je dis que $Y'$ est une partie fermée de \deleted{\unsure{N}}$M$ ; \deleted{et que si \inserted{[\unsure{cela se}]}} on munit $Y'$ de la \deleted{\unsure{topol.}} structure réduite induite [NB Alors $Y'\to Y$ est un morphisme \deleted{\unsure{fini}} \inserted{\unsure{propre} radiciel \deleted{\unsure{sur}}\unsure{surjectif}, donc fini...} \deleted{[Cela \unsure{fait|soit} admis}, On aura terminé, car \deleted{\unsure{si on|$\varphi$}} on aura un $Y$-morphisme $Y'\to M$, correspondant à un sous-schéma \inserted{fermé} \deleted{X} $Z$ de $X'$, plat de prés finie sur $Y$, dont les fibres sont les $(X_y)_{\mathrm{réd}}\otimes_{k(y)}k(y')$, \unsure{ens.|donc|qui sont} géom. \unsure{normales|réduites}.. Par suite, $Z$ est \unsure{normal|réduit} sur $Y'$ \inserted{réduit,} donc réduit, et ayant \unsure{même} ensemble sous-jacent que $X'$, il est égal à $X'_{\mathrm{réd}}$].

% archive sheet number: 11

\scanpage{13}
\diagram{\begin{tikzcd}[ampersand replacement=\&]
\unsure{G|\mathcal G}\arrow[r,"u"] \& F\arrow[d]\\
\& \deleted{\unsure{G|S}}S\ \text{loc. noeth}
\end{tikzcd}}

\deleted{Si \unsure{$G$ et $F$} \unsure{est}}\inserted{Supposons}

(i) $F$ et $G$ sont des faisceaux fpqc, commutent aux $\unsure{\underline{\lim}|\underset{\rightarrow}{\lim}}$ d'anneaux, aux $\underset{\leftarrow}{\lim}$ adiques d'anneaux artiniens

(ii) $u$ est un monomorphisme, et est infinitésimalement étale.

\deleted{(iii) $u$ est un \unsure{épi} \unsure{épimorphisme|effectif|épimorphe}}

Alors $G\unsure{\mathbin{-}|\to}F$ est une immersion ouverte.

\noindent\rule{.35\linewidth}{.4pt}

Une Immersion \unsure{ouverte} surjective de foncteurs est un isomorphisme

\noindent\rule{.39\linewidth}{.4pt}

Caractérisation des morphismes étales de foncteurs...,

% archive sheet number: 12

\scanpage{14}
\figure{\resizebox{\linewidth}{!}{%
\begin{tikzpicture}[x=.010cm,y=-.010cm,every node/.style={inner sep=1pt,font=\small},line cap=round,line join=round]
\node[anchor=west] at (616,102) {$X\to Y$ une imm ouverte};
\node[anchor=west] at (228,156) {\parbox{2.1cm}{\diagram{\begin{tikzcd}[ampersand replacement=\&,column sep=small] X\arrow[r] \& Y \end{tikzcd}}}};
\node at (270,414) {$S$};

\draw (587,220) -- (1539,146) .. controls (1591,139) and (1592,163) .. (1590,196) -- (1585,396);
\draw (597,233) -- (617,497) -- (1544,427);
\node[anchor=west] at (656,269) {morphisme \deleted{\unsure{ét|de|affine}} $\mathrm{form}^{\mathrm t}$ lisse};
\node[anchor=west] at (1274,258) {$\mathrm{inf}^{\mathrm t}$ lisse};
\node[anchor=west] at (936,324) {--- non ram};
\node[anchor=west] at (1274,314) {$\mathrm{inf}^{\mathrm t}$ \deleted{\unsure{étale}}\unsure{non ram.}};
\node[anchor=west] at (936,374) {--- étale};
\node[anchor=west] at (1274,371) {$\mathrm{inf}^{\mathrm t}$ étale};

\draw (677,483) .. controls (723,465) and (928,461) .. (1188,435) .. controls (1212,431) and (1212,439) .. (1217,454) -- (1237,645) -- (736,686) -- (693,506);
\node[anchor=west] at (737,507) {immersion};
\node[anchor=west] at (743,571) {\deleted{immersion \uline{ouverte}}};
\node[anchor=west] at (754,639) {immersion \uline{fermée}};
\draw (713,564) -- (541,578) -- (550,640) -- (723,635);
\node at (629,605) {\unsure{Mn|Mon}};
\draw (1237,506) -- (1588,482) .. controls (1602,479) and (1608,495) .. (1609,557) -- (1240,601);
\node at (1424,537) {imm ouverte};

\draw (344,686) -- (344,881) .. controls (359,902) and (596,878) .. (627,871) .. controls (651,865) and (644,798) .. (635,691) .. controls (635,656) and (610,654) .. (571,663) -- (353,677);
\draw (395,724) -- (379,724) -- (392,869) -- (405,869);
\node[anchor=west] at (404,754) {séparé};
\node at (459,815) {$\sim$};
\draw (684,709) -- (956,698) -- (973,893);
\draw (691,716) -- (713,929) -- (928,905) .. controls (947,902) and (962,898) .. (964,893);
\draw (746,715) -- (734,713) .. controls (725,728) and (750,814) .. (739,868) -- (756,871);
\node[anchor=west] at (763,728) {propre};
\foreach \yy in {785,801,817,841} \fill (824,\yy) circle[radius=.9pt];

\draw (1025,770) -- (1304,750) .. controls (1321,756) and (1328,806) .. (1332,878);
\draw (1030,793) -- (1032,902) -- (1315,883);
\node at (1171,844) {faisceau};
\draw (1018,935) -- (1485,889) -- (1494,1068) -- (1020,1106) -- (1024,958);
\node at (1222,966) {\uline{loc. de prés. finie}};
\node at (1281,1045) {[\uline{relatif}]};
\draw (1025,1133) -- (1384,1117) -- (1395,1230);
\draw (1028,1140) -- (1034,1228) -- (1411,1203);
\node at (1145,1181) {adiques};
\draw (1009,786) .. controls (998,831) and (1009,925) .. (954,979) .. controls (999,986) and (981,1034) .. (960,1158) .. controls (955,1183) and (979,1194) .. (990,1204);
\node at (945,981) {\deleted{\unsure{les|Ex.}}};
\node[draw,circle,inner sep=3pt] at (865,1078) {$A$};

\node at (334,943) {\parbox{2.1cm}{\diagram{\begin{tikzcd}[ampersand replacement=\&,column sep=small] F_i\arrow[r] \& F \end{tikzcd}}}};
\node[align=center] at (582,966) {\uline{imm}\\\uline{ouverte}};
\node at (328,1116) {\parbox{2.25cm}{\diagram{\begin{tikzcd}[ampersand replacement=\&,column sep=small,row sep=small]
\deleted{\unsure{S|G}}G\arrow[r] \& F\\
G_T\arrow[u]\arrow[r] \& \deleted{\unsure{S}}T\arrow[u,bend right=10]
\end{tikzcd}}}};

\draw (448,1230) -- (763,1216) .. controls (781,1218) and (779,1253) .. (787,1398) -- (481,1402);
\draw (444,1251) -- (459,1401);
\node at (610,1273) {connexe};
\node at (624,1350) {$\vdots$};
\end{tikzpicture}}}

% archive sheet number: 13
